\documentclass{article}

\usepackage{changepage}
\usepackage{xcolor}
\usepackage{amsmath}
\usepackage{amssymb}
\usepackage{tikz}
\usepackage{tcolorbox}

\definecolor{blue}{RGB}{0, 166, 223}
\definecolor{lightblue}{RGB}{220, 240, 247}

\newcommand{\definition}[1]{
	\vspace{0.25cm}
	\begin{tcolorbox}[
		colback=lightblue,
		colframe=blue,
		colbacktitle=blue,
		coltitle=white,
		title=\textbf{Definition},
		boxrule=0.5pt,
		arc=0pt,
		left=0.4cm,
		right=0.4cm,
		top=0.25cm,
		bottom=0.35cm,
		toptitle=0.1cm,
		bottomtitle=0.075cm
	]
		#1
	\end{tcolorbox}
	\vspace{0.25cm}
}

\newcommand{\theorem}[2]{
	\vspace{0.25cm}
	\begin{tcolorbox}[
		colback=lightblue,
		colframe=blue,
		boxrule=0.5pt,
		toprule=2pt,
		arc=0pt,
		left=0.4cm,
		right=0.4cm,
		top=0.25cm,
		bottom=0.35cm
	]
		\textcolor{blue}{\textbf{Theorem #1}}

		\vspace{0.15cm}

		#2
	\end{tcolorbox}
	\vspace{0.25cm}
}

\newcommand{\container}[1]{
	\vspace{0.25cm}
		\begin{adjustwidth}{0cm}{0cm}
			#1
		\end{adjustwidth}
	\vspace{0.25cm}
}

\newcommand{\padding}[1]{
	\vspace{0.25cm}
		\begin{adjustwidth}{0cm}{0cm}
			#1
		\end{adjustwidth}
	\vspace{0cm}
}

\newcommand{\calc}[1]{
	\[
		#1
	\]
}

\begin{document}

	\container{
		\textcolor{blue}{\textbf{5.1.11 Properties of Summations and Products}}
	}

	\container{
		The following theorem states general properties 
		of summations and products. The proof of the theorem
		is discussed in Section 5.6.
	}

	\theorem{5.1.1}{
		If \(a_m, a_{m+1}, a_{m+2}, \text{...}\) and \(b_m, b_{m+1},
		b_{m+2}, \text{...}\) are sequences of real numbers and c is
		any real number, then the following equations hold for any
		integer \(n \geq m\):

		\calc{
			\begin{aligned}
				&\text{1.} \sum_{k=m}^{n} a_k + \sum_{k=m}^{n} b_k = \sum_{k=m}^{n} \left(a_k + b_k\right) \\
				&\text{2. } c \times \sum_{k=m}^{n} a_k = \sum_{k=m}^{n} c \times a_k \textcolor{blue}{\tiny\text{ generalized distributive law}} \\
				&\text{3.} \left(\prod_{k=m}^{n} a_k\right) \times \left(\prod_{k=m}^{n} b_k\right) = \prod_{k=m}^{n} \left(a_k \times b_k\right)
			\end{aligned}
		}
	}

\end{document}
